% Vietnamese translation for OI Public Library
% Source-PDF: IOI中国国家候选队论文/2004/栗师.pdf
\documentclass[11pt]{article}
\usepackage{oipl-vn}

\newcommand{\Z}{\mathbb{Z}}
\newcommand{\N}{\mathbb{N}}
\newcommand{\lcm}{\operatorname{lcm}}
\newcommand{\sgn}{\operatorname{sgn}}

\title{Ứng dụng chuyển hóa mục tiêu trong giải bài}
\author{Li Shi\\Trường Trung học Changjun, Trường Sa, Hồ Nam}
\date{IOI 2004, đội tuyển tập huấn quốc gia Trung Quốc}

\begin{document}
\maketitle

\origtitle{Ứng dụng của chuyển hóa mục tiêu trong giải bài}
\sourcepath{IOI China candidate team papers / 2004 / Li Shi.pdf}

\begin{abstract}
Bài viết thảo luận ngắn gọn việc ứng dụng tư tưởng chuyển hóa mục tiêu trong
thiết kế thuật toán và trong phân tích bài toán.

Phần thứ nhất nêu vì sao cần chuyển hóa mục tiêu. Phần thứ hai dùng ví dụ để
giải thích vai trò của chuyển hóa mục tiêu trong thiết kế thuật toán: trước
hết đạt tới mục tiêu đã được chuyển hóa, rồi từ đó suy ra mục tiêu cuối cùng.
Phần thứ ba cũng thông qua một ví dụ để giới thiệu tác dụng của chuyển hóa mục
tiêu trong phân tích bài toán. Cuối cùng, bài viết tổng kết một số cách
chuyển hóa mục tiêu thường gặp và bàn cách vận dụng chúng linh hoạt.
\end{abstract}

\noindent\textbf{Từ khóa:} chuyển hóa mục tiêu; phóng đại; thu hẹp; đơn giản
hóa bài toán.

\section{Mở đầu}

Trong quá trình thiết kế thuật toán tin học, ta luôn gặp khó khăn kiểu này hay
kiểu khác. Một nguyên nhân lớn là chiều sâu và chiều rộng tư duy của con
người đều hữu hạn, trong khi thế giới chưa biết là vô hạn. Khi gặp một bài
toán khó giải, ta thường cảm thấy mục tiêu quá xa, quan hệ chằng chịt, không
biết bắt đầu từ đâu. Khi ấy, có thể thử chuyển hóa mục tiêu và suy nghĩ trên
mục tiêu mới: giảm bớt ràng buộc, nới lỏng điều kiện, thu hẹp phạm vi, v.v.
Sau khi giải được mục tiêu đã chuyển hóa, ta đứng ở một độ cao mới để thiết
kế thuật toán, hoặc mở rộng thuật toán của mục tiêu mới về mục tiêu ban đầu.
Tư tưởng chuyển hóa mục tiêu cung cấp cho ta lối tắt theo hai hướng: thuật
toán và đường lối suy nghĩ.

\section{Chuyển hóa mục tiêu trong thiết kế thuật toán}

Nếu một bước chưa đạt được đích, ta chia thành nhiều bước; mỗi bước đạt một
mục tiêu trung gian. Cách này được dùng rất thường xuyên. Dạng thường gặp nhất
là: nếu mục tiêu có nhiều ràng buộc, trước hết chỉ xét một ràng buộc để rút ra
kết luận, rồi từ kết luận ấy xét các ràng buộc còn lại. Dưới đây dùng một bài
của Polish Olympiad in Informatics 2003 để minh họa phương pháp này.

\subsection*{Ví dụ 1: Siêu mã}

Trên một bàn cờ vô hạn có một quân siêu mã. Nó có thể thực hiện nhiều loại
nước đi. Mỗi loại nước đi được xác định bởi hai số nguyên: số thứ nhất cho
biết số cột dịch chuyển, dương là sang phải và âm là sang trái; số thứ hai cho
biết số hàng dịch chuyển, dương là lên trên và âm là xuống dưới.

\paragraph{Nhiệm vụ.}
Viết chương trình đọc từ tệp \texttt{SUP.IN} cơ sở dữ liệu mô tả nhiều quân
siêu mã. Với mỗi quân, hãy xác định nó có thể tự đi tới mọi ô trên bàn cờ hay
không. Ghi kết quả ra tệp \texttt{SUP.OUT}.

\paragraph{Dữ liệu vào.}
Dòng đầu của \texttt{SUP.IN} chứa số nguyên \(K\), là số cơ sở dữ liệu,
\(1\le K\le 100\). Sau đó là \(K\) bộ dữ liệu. Dòng đầu mỗi bộ chứa số nguyên
\(N\), là số loại nước đi của quân mã, \(1\le N\le 100\). Mỗi dòng trong
\(N\) dòng tiếp theo chứa hai số nguyên \(P,Q\), cách nhau bởi một khoảng
trắng, mô tả một nước đi; \(-100\le P,Q\le 100\).

\paragraph{Dữ liệu ra.}
Tệp \texttt{SUP.OUT} gồm \(K\) dòng. Nếu quân siêu mã ở bộ dữ liệu thứ \(i\)
có thể tới mọi ô trên bàn cờ thì dòng thứ \(i\) chứa \texttt{TAK}; ngược lại
chứa \texttt{NIE}.

\begin{center}
\begin{tabular}{ll}
\texttt{SUP.IN} & \texttt{SUP.OUT}\\
\begin{minipage}{0.35\linewidth}
\begin{verbatim}
2
5
2 4
2 2
-3 -3
4 3
1 3
5
1 -3
2 1
4 1
4 -2
2 -2
\end{verbatim}
\end{minipage}
&
\begin{minipage}{0.2\linewidth}
\begin{verbatim}
TAK
NIE
\end{verbatim}
\end{minipage}
\end{tabular}
\end{center}

\subsection{Xác định mô hình thuật toán}

Nhìn đề bài, thuật toán dễ nghĩ nhất là tìm kiếm theo chiều rộng. Từ những ô
đã biết là đi tới được, ta dùng các nước đi của quân mã để mở rộng ra các ô
khác, cứ thế tiếp tục, cuối cùng phán đoán có mở rộng hết cả bàn cờ hay
không.

Ta nhanh chóng thấy cách này không được: số ô của bàn cờ là vô hạn, nên không
thể kiểm tra đã tới được tất cả các ô. Khó khăn này có vẻ được giải quyết
ngay: hiển nhiên chỉ cần phán đoán siêu mã có thể tới bốn ô kề điểm xuất phát
hay không. Nếu tới được bốn ô ấy, nó tất nhiên tới được mọi vị trí trên bàn cờ.

Vấn đề đã xong chưa? Chưa. Tuy mục tiêu cuối cùng chỉ cần xét bốn ô, quá trình
đi tới bốn ô ấy có thể đi qua rất nhiều điểm. Tệ hơn, khi bài toán vô nghiệm
ta không có cách phán đoán đúng; lúc cài đặt thuật toán dễ rơi vào vòng lặp
không dừng.

Một ý khác là biến bàn cờ vô hạn thành một bàn cờ hữu hạn đủ lớn, rồi xem trên
bàn cờ hữu hạn ấy có tới được bốn ô kia hay không. Nhưng đây vẫn là việc vô
ích: bàn cờ hữu hạn phải rất lớn, hiệu suất thời gian thấp và thuật toán cũng
thiếu chứng minh.

Thử các thuật toán đồ thị, quy hoạch động, tham lam, v.v., cuối cùng đều đi
tới kết quả có thể đoán trước: thất bại. Muốn có thuật toán hiệu quả, dường
như chỉ còn một con đường: dùng tư tưởng toán học.

Muốn giải bằng toán học, trước hết phải xây dựng mô hình toán học. Lấy ô ban
đầu của siêu mã làm gốc tọa độ trong mặt phẳng. Một loại nước đi của quân mã
được biểu diễn bằng vectơ phẳng \(P_i=(x_i,y_i)\). Khi đó ta cần phán đoán:
với mọi \(x,y\), có tồn tại một bộ
\[
  c_1,c_2,\ldots,c_n,\qquad c_i\ge 0,\quad 1\le i\le n,
\]
sao cho
\[
  \sum_{i=1}^{n} c_iP_i=(x,y)
\]
hay không.

Như vậy, mô hình toán học của bài toán là giải phương trình.

Hơn nữa, chỉ cần kiểm tra bốn trường hợp
\[
  (x,y)=(-1,0),(1,0),(0,-1),(0,1).
\]
Bốn trường hợp ấy về bản chất giống nhau, nên ta chỉ xét \((0,1)\). Cần phán
đoán phương trình sau có nghiệm nguyên không âm hay không:
\[
  c_1P_1+c_2P_2+\cdots+c_nP_n=(0,1). \tag{1}
\]

Ví dụ, với bộ dữ liệu thứ nhất trong đề: \(n=5\) và
\[
P_1=(2,4),\quad P_2=(2,2),\quad P_3=(-3,-3),\quad
P_4=(4,3),\quad P_5=(1,3),
\]
một nghiệm nguyên không âm của (1) là
\[
  3(2,4)+5(2,2)+13(-3,-3)+5(4,3)+3(1,3)=(0,1).
\]

\subsection{``Phóng đại'' mục tiêu}

Đây là một phương trình tuyến tính, số ẩn tương đối nhiều trong khi chỉ có một
phương trình vectơ. Trực giác cho thấy nếu phương trình có nghiệm thì số
nghiệm hẳn rất nhiều, rất có thể là vô hạn. Vì vậy suy nghĩ theo hướng dựng
một nghiệm là lựa chọn hợp lý hơn.

Sau khi thử nhiều cách dựng khác nhau, ta thấy có hai yếu tố cản trở: một là
phải làm cho tung độ vế phải bằng \(1\); hai là mọi hệ số \(c\) ở vế trái phải
không âm. Trước mặt ta là hai khối đá nặng. Khuân cả hai cùng lúc thì quá
nặng; vậy hãy khuân từng khối. Trước hết dựng một nghiệm sao cho vế phải của
(1) bằng \((0,1)\), còn mỗi hệ số \(c\) ở vế trái chỉ cần là số nguyên.

Vậy ta xác định một mục tiêu đã phóng đại: tìm nghiệm nguyên của phương trình.

Tìm nghiệm nguyên vẫn còn khó. Ta bắt đầu từ trường hợp đơn giản
\(n=2\), \(P_1=(x_1,y_1)\), \(P_2=(x_2,y_2)\). Để trình bày gọn, giả sử
\(x_1,y_1,x_2,y_2\ne 0\). Siêu mã có tới được \((0,1)\) hay không?
Trước hết, hai hệ số \(c_1,c_2\) phải thỏa
\[
  c_1x_1+c_2x_2=0.
\]
Gọi \(m>0\) là bội chung nhỏ nhất của \(|x_1|\) và \(|x_2|\). Đặt
\[
  c_1=\frac{m}{x_1},\qquad c_2=-\frac{m}{x_2}.
\]
Khi đó
\[
  W=c_1P_1+c_2P_2
\]
là một vectơ thẳng đứng sinh bởi \(P_1,P_2\), và là vectơ có độ dài nhỏ nhất
theo nghĩa này. Nếu \(W\) có thể được dựng từ \(P_1,P_2\), tức tồn tại
\(c_1,c_2\in\Z\) sao cho \(c_1P_1+c_2P_2=W\), thì với mọi \(k\in\Z\), vectơ
\(kW\) cũng dựng được từ \(P_1,P_2\). Nếu \(|W|=1\), trường hợp hai vectơ có
nghiệm nguyên; ngược lại, nếu \(|W|\ne 1\), không có nghiệm nguyên.

Sau khi giải trường hợp hai vectơ, quay lại trường hợp \(n\) vectơ, ta tự
nhiên nảy ra ý tưởng: dùng các vectơ thẳng đứng thu được từ từng cặp vectơ để
dựng các vectơ thẳng đứng khác.

Ở trên ta dùng \(P_1,P_2\) dựng được dãy vectơ thẳng đứng \(kW\). Gọi tập các
tung độ của những vectơ ấy là \(S_{1,2}\). Với mọi cặp khác nhau
\(P_i,P_j\), \(i<j\), ta đều có thể dựng một tập \(S_{i,j}\).

Vẫn xét bộ dữ liệu thứ nhất:
\[
P_1=(2,4),\ P_2=(2,2),\ P_3=(-3,-3),\ P_4=(4,3),\ P_5=(1,3).
\]
Ta thu được bảng sau:
\[
\begin{array}{c|cccc}
 &2&3&4&5\\ \hline
1&\{2k\mid k\in\Z\}&\{6k\mid k\in\Z\}&\{5k\mid k\in\Z\}&\{2k\mid k\in\Z\}\\
2&&\{0\}&\{k\mid k\in\Z\}&\{4k\mid k\in\Z\}\\
3&&&\{3k\mid k\in\Z\}&\{6k\mid k\in\Z\}\\
4&&&&\{9k\mid k\in\Z\}
\end{array}
\]

Gộp tất cả \(S_{i,j}\), xem siêu mã còn tới được những vị trí nào trên trục
tung. Định nghĩa tập tung độ \(S\) như sau:
\begin{enumerate}
  \item Nếu tồn tại \(i,j\) sao cho \(Y\in S_{i,j}\), thì \(Y\in S\).
  \item Nếu \(Y_1,Y_2\in S\) và \(Y=k_1Y_1+k_2Y_2\) với \(k_1,k_2\in\Z\), thì
  \(Y\in S\).
  \item Ngoài các giá trị được định nghĩa bởi hai quy tắc trên, không có
  \(Y\) nào thuộc \(S\).
\end{enumerate}
Thực chất, \(S\) là tập các số có thể nhận được từ các phần tử trong
\(S_{i,j}\) bằng phép cộng và trừ.

Theo kiến thức về số dư, \(S\) cũng có dạng
\[
  \{kY\mid k\in\Z\},
\]
trong đó \(Y\) là ước chung dương lớn nhất của tất cả các số trong các
\(S_{i,j}\). Nếu \(Y=1\), phương trình (1) có nghiệm nguyên. Với bảng trên,
\[
  S=\{k\mid k\in\Z\},
\]
nên có thể dựng được một nghiệm nguyên. Chẳng hạn lấy
\[
S_{1,2}=\{2k\},\qquad S_{4,5}=\{9k\}.
\]
Vì ước chung lớn nhất của \(2\) và \(9\) là \(1\), và \(-4\cdot2+9=1\), ta có
\[
  -4(P_1-P_2)+(4P_5-P_4)=-4P_1+4P_2-P_4+4P_5=(0,1).
\]

Như vậy ta đã giải được một nửa bài toán tìm nghiệm nguyên: tìm ra một điều
kiện đủ để (1) có nghiệm nguyên.

Nhưng cách dựng trên chỉ là điều kiện đủ. Nó có cần thiết không? Nói cách
khác, nếu một bộ \(P\) làm cho (1) có nghiệm nguyên
\(c_1,c_2,\ldots,c_n\), thì liệu \((0,1)\) nhất định có thể tạo ra từ một số
vectơ thẳng đứng có tung độ thuộc các \(S_{i,j}\) bằng phép cộng trừ hay
không? Câu trả lời là có.

\paragraph{Kết luận 2.2.1.}
Có thể chia vế trái của (1), một cách tương đương, thành tổng của một số đa
thức; mỗi đa thức có nhiều nhất hai hạng tử, và tổng của mỗi đa thức là một
vectơ thẳng đứng.

Ví dụ,
\[
  3(2,4)-13(2,2)+9(-3,-3)+11(4,3)+3(1,3)=(0,1)
\]
có thể biến thành
\[
\begin{aligned}
&[3(2,4)-3(2,2)]
+[-10(2,2)+5(4,3)]\\
&\quad +[9(-3,-3)+27(1,3)]
+[6(4,3)-24(1,3)]\\
&=(0,6)+(0,-5)+(0,54)+(0,-54)=(0,1).
\end{aligned}
\]

Chứng minh kết luận bằng quy nạp. Nếu \(n\le 2\), vế trái tự nó đã là một đa
thức có nhiều nhất hai hạng tử, nên hiển nhiên đúng. Giả sử kết luận đúng khi
\(n=k-1\). Ta chứng minh cho \(n=k\).

Trong chứng minh chỉ cần quan tâm tới phần hoành độ của các \(P\), không cần
quan tâm tới phần tung độ. Vì
\[
  \sum_{i=1}^{k} c_ix_i=0,
\]
nếu tồn tại hai dãy số nguyên \(u_i,v_i\) sao cho biểu thức trên có thể biến
đổi tương đương thành
\[
  \sum_{i=2}^{k}(u_ix_i+v_ix_1)+
  \sum_{i=2}^{k}(c_i-u_i)x_i=0,
\]
trong đó
\[
  u_ix_i+v_ix_1=0,\qquad \sum_{i=2}^{k}v_i=c_1,
\]
thì ta có thể lấy ra từ vế trái \(k-1\) đa thức \(u_ix_i+v_ix_1\). Mỗi đa
thức ấy có không quá hai hạng tử và có tổng hoành độ bằng \(0\).

Có thể giả sử ước chung lớn nhất của mọi \(x_i\) là \(1\); nếu không, chia tất
cả \(x_i\) cho ước chung ấy, điều này không ảnh hưởng tới việc chọn \(u,v\).
Đặt
\[
  g_i=\frac{\lcm(x_i,x_1)}{x_1}\qquad (i\ge 2).
\]
Không khó thấy \(g_i\mid x_i\). Gọi
\[
  g=\gcd(g_2,g_3,\ldots,g_k),
\]
thì \(g\mid x_i\) với \(i\ge 2\). Từ giả thiết ước chung lớn nhất của mọi
\(x_i\) bằng \(1\), suy ra \(g\) nguyên tố cùng nhau với \(x_1\). Lại vì
\[
  g\mid \sum_{i=2}^{k}c_ix_i,
\]
nên
\[
  g\mid c_1x_1,
\]
tức \(g\mid c_1\). Do đó tồn tại dãy số nguyên \(d_2,d_3,\ldots,d_k\) sao cho
\[
  \sum_{i=2}^{k}d_ig_i=c_1.
\]
Đặt
\[
  v_i=d_ig_i,\qquad
  u_i=-\frac{v_ix_1}{x_i}
     =-\frac{d_ig_ix_1}{x_i}
     =-\frac{d_i\lcm(x_1,x_i)}{x_i}.
\]
Ta thu được các dãy \(u,v\) cần thiết. Vậy kết luận 2.2.1 được chứng minh.

Qua từng bước tính toán và chứng minh, ta hoàn thành bước đầu của bài toán:
tìm được điều kiện cần và đủ để (1) có nghiệm nguyên.

\paragraph{Kết luận 2.2.2.}
Phương trình (1) có nghiệm nguyên khi và chỉ khi \((0,1)\) có thể nhận được
bằng phép cộng trừ từ các vectơ thẳng đứng \(W\) dạng sau: tồn tại
\[
  i,j,k_1,k_2,\qquad 1\le i<j\le n,\quad k_1,k_2\in\Z,
\]
sao cho
\[
  k_1P_i+k_2P_j=W.
\]

\subsection{Nhắm lại mục tiêu}

Nếu ba vectơ còn lại \((0,-1),(-1,0),(1,0)\) cũng được kiểm tra tương tự, ta
có thể phán đoán phương trình có nghiệm nguyên hay không. Phần kết luận tiếp
theo có tiền đề lớn là phương trình đã có một nghiệm nguyên. Để trình bày và
chứng minh gọn hơn, ta còn có thể giả sử trong các vectơ đã cho có một vectơ
có hoành độ dương, một vectơ có hoành độ âm; tung độ cũng vậy.

Vì tư tưởng tìm nghiệm nguyên đã cho hiệu quả tốt, ta tiếp tục dùng nó để tìm
nghiệm nguyên không âm. Tổng hợp lại, phần trên dùng hai kỹ thuật: bắt đầu từ
hai vectơ, và trước hết dựng một điều kiện đủ cho việc có nghiệm rồi chứng
minh điều kiện đủ ấy là cần thiết.

Ta lại bắt đầu từ hai vectơ. Lúc này vì yêu cầu nghiệm nguyên không âm, phạm
vi từ số nguyên chuyển thành số nguyên không âm.

Theo cách định nghĩa \(S_{i,j}\), ta định nghĩa \(T_{i,j}\): nếu một vectơ
thẳng đứng \(W\) có thể biểu diễn dưới dạng
\[
  k_1P_i+k_2P_j,\qquad k_1,k_2\in\N,
\]
thì tung độ của \(W\) thuộc \(T_{i,j}\). Hiển nhiên, nếu \(P_i,P_j\) cùng
hướng theo trục \(x\), thì \(T_{i,j}\) là rỗng.

Vẫn với ví dụ
\[
P_1=(2,4),\ P_2=(2,2),\ P_3=(-3,-3),\ P_4=(4,3),\ P_5=(1,3),
\]
ta được bảng:
\[
\begin{array}{c|cccc}
 &2&3&4&5\\ \hline
1&\varnothing&\{6k\mid k\in\N\}&\varnothing&\varnothing\\
2&&\{0\}&\varnothing&\varnothing\\
3&&&\{-3k\mid k\in\N\}&\{6k\mid k\in\N\}\\
4&&&&\varnothing
\end{array}
\]

Định nghĩa tập \(T\):
\begin{enumerate}
  \item Nếu tồn tại \(i,j\) sao cho \(Y\in T_{i,j}\), thì \(Y\in T\).
  \item Nếu \(Y_1,Y_2\in T\) và \(Y=k_1Y_1+k_2Y_2\) với \(k_1,k_2\) là số
  nguyên không âm, thì \(Y\in T\).
  \item Ngoài các giá trị được định nghĩa bởi hai quy tắc trên, không có
  \(Y\) nào thuộc \(T\).
\end{enumerate}
Không khó thấy \(T\) là tập tung độ của các vectơ thẳng đứng nhận được từ
các \(T_{i,j}\) bằng phép cộng. Ví dụ trên cho
\[
  T=\{3k\mid k\in\Z\}.
\]

Vậy có phải chỉ cần \((0,1)\in T\) là đủ không? Rõ ràng không; nếu không thì
quá trình tìm nghiệm nguyên ở trên còn có ích gì? Phản ví dụ rất dễ tìm:
\[
  P_1=(3,0),\quad P_2=(-1,1),\quad P_3=(-1,-1).
\]
Quân siêu mã có thể tới mọi điểm, nhưng trong \(T\) chỉ có các vectơ thẳng
đứng có độ dài là bội của \(3\). Vì vậy cần dùng kết luận về nghiệm nguyên ở
trên.

Hiện tại ta đã có một nghiệm nguyên của phương trình. Làm thế nào tìm một
nghiệm nguyên không âm? Một ý tưởng là điều chỉnh: liên tục cộng vectơ không
vào vế trái của phương trình để hệ số mỗi hạng tử đều tăng lên; hoặc trước
hết làm cho các hệ số bên trái đều dương, rồi dùng nghiệm nguyên đã tìm được
để biến vế phải thành \(1\). Từ cách dựng ấy có thể rút ra kết luận sau.

\paragraph{Kết luận 2.3.1.}
Nếu \(T\) có một số dương và một số âm, thì phương trình (1) nhất định có
nghiệm nguyên không âm.

\paragraph{Chứng minh.}
Chỉ cần dựng một nghiệm nguyên không âm. Vì trong các vectơ đã cho có một
vectơ có hoành độ dương và một vectơ có hoành độ âm, tồn tại một dãy số nguyên
dương
\[
  d_1,d_2,\ldots,d_n
\]
rất lớn sao cho
\[
  \sum_{i=1}^{n}d_ix_i=0.
\]
Đặt
\[
  M=\sum_{i=1}^{n}d_iy_i.
\]
Do \(T\) có một số dương và một số âm, ta có thể tìm trong \(T\) một số sao
cho khi cộng vào \(M\) thì giá trị thu được gần \(0\) nhất; cộng các hệ số
tương ứng của vectơ ấy vào vế phải. Sau đó dùng nghiệm nguyên của \((0,1)\)
hoặc \((0,-1)\) đã tìm được, cộng hệ số của nó vào tiếp để biến \(M\) thành
\(1\). Vì các \(d_i\) ban đầu được chọn rất lớn, các hệ số âm cộng thêm sau
đó nhỏ so với chúng, nên mọi hệ số cuối cùng vẫn là số nguyên không âm. Do đó
phương trình nhất định có nghiệm nguyên không âm.

Ta thực hiện các bước trên cho bộ dữ liệu thứ nhất. Chọn dãy
\[
  d=(9,9,27,9,9),
\]
vì
\[
  9\cdot2+9\cdot2+27\cdot(-3)+9\cdot4+9\cdot1=0.
\]
Khi đó
\[
  M=9y_1+9y_2+27y_3+9y_4+9y_5=27.
\]
Chọn số âm trong \(T\) là \(-27\). Vì
\[
  (0,-27)=9(4P_3+3P_4)=36P_3+27P_4,
\]
cộng \(36y_3+27y_4\) vào \(M\), ta được
\[
  M=9y_1+9y_2+63y_3+36y_4+9y_5=0.
\]
Sau đó dùng nghiệm nguyên đã có
\[
  -4P_1+4P_2-P_4+4P_5=(0,1),
\]
tức cộng
\[
  -4y_1+4y_2-y_4+4y_5=1
\]
vào \(M\), ta được
\[
  M=5y_1+13y_2+63y_3+35y_4+13y_5=1.
\]
Nói cách khác,
\[
  5P_1+13P_2+63P_3+35P_4+13P_5=(0,1)
\]
là một nghiệm nguyên không âm cần tìm.

Tiếp theo chỉ cần chứng minh điều kiện còn lại.

\paragraph{Kết luận 2.3.2.}
Nếu phương trình (1) có nghiệm nguyên không âm, thì tồn tại
\[
  P_i,P_j,k_1,k_2,\qquad k_1,k_2\ge 0,
\]
sao cho \(k_1P_i+k_2P_j\) là một vectơ nằm trên nửa trục dương của trục
\(Y\).

\paragraph{Chứng minh.}
Để trình bày gọn, giả sử trong các vectơ không có vectơ nằm ngang hoặc thẳng
đứng, tức không có vectơ nào có \(x=0\) hoặc \(y=0\). Chứng minh bằng phản
chứng. Giả sử không tồn tại \(P_i,P_j,k_1,k_2\) như vậy. Khi đó trong các
vectơ \(P\) không thể đồng thời có vectơ dạng \((a,b)\) và dạng \((-c,d)\),
với \(a,b,c,d\) đều là số nguyên dương, vì hai vectơ này chắc chắn dựng được
một vectơ trên nửa trục dương của \(Y\).

Không mất tính tổng quát, giả sử các vectơ có tung độ dương chỉ có dạng
\((a,b)\); gọi tập các vectơ này là \(Q_1\). Còn có một số vectơ dạng
\((-a,-b)\); gọi tập này là \(Q_2\). Những vectơ còn lại thuộc \(Q_3\), tức có
dạng \((a,-b)\). Có thể thấy \(Q_1,Q_2\) đều không rỗng. Khi ấy nhất định có
\[
  t_1=\max\left\{\frac{y_i}{x_i}\mid (x_i,y_i)\in Q_1\right\}
  <
  \min\left\{\frac{y_i}{x_i}\mid (x_i,y_i)\in Q_2\right\}=t_2.
\]
Trong hình gốc, hai đường thẳng có hệ số góc \(t_1,t_2\) kẹp các nhóm
\(Q_1,Q_2,Q_3\), trực quan cho thấy không thể đi tới nửa trục dương của
\(Y\).

Về đại số, từ một nghiệm không âm \(\sum c_iP_i=(0,1)\), phần hoành độ bằng
\(0\). Suy ra
\[
\begin{aligned}
0
&\ge -\sum_{P_i\in Q_3}c_ix_i
 = \sum_{P_i\in Q_1}c_ix_i+\sum_{P_i\in Q_2}c_ix_i\\
&\ge
  \sum_{P_i\in Q_1} c_ix_i\frac{y_i/x_i}{t_1}
 +\sum_{P_i\in Q_2} c_ix_i\frac{y_i/x_i}{t_2}\\
&>
  \sum_{P_i\in Q_1} \frac{c_iy_i}{t_2}
 +\sum_{P_i\in Q_2} \frac{c_iy_i}{t_2}
 =
  \frac{\sum_{P_i\in Q_1\cup Q_2}c_iy_i}{t_2}\\
&=
  -\frac{\sum_{P_i\in Q_3}c_iy_i}{t_2}.
\end{aligned}
\]
Vì \(t_2>0\), suy ra
\[
  \sum_{P_i\in Q_3}c_iy_i>0,
\]
mâu thuẫn với việc mọi vectơ trong \(Q_3\) đều có dạng \((a,-b)\). Do đó kết
luận 2.3.2 đúng.

Từ đó thu được thuật toán cuối cùng của bài toán. Dưới đây là giả mã chính.

\begin{verbatim}
Buoc 1:
g <- 0
For moi Pi la vector thang dung:
    g <- gcd(g, |yi|)
For moi Pi khong thang dung:
    For moi Pj khong thang dung:
        t <- gcd(|xi|, |xj|)
        len <- |yi*xj/t - yj*xi/t|
        g <- gcd(g, len)

Buoc 2:
If g != 1:
    in "NIE"; dung
positive <- false
negative <- false
For moi Pi la vector thang dung:
    If yi > 0 then positive <- true
    Else negative <- true
For moi Pi khong thang dung:
    For moi Pj khong thang dung va xi, xj trai dau:
        W <- Pi*|xj| + Pj*|xi|
        If W la vector huong duong theo truc Y:
            positive <- true
        If W la vector huong am theo truc Y:
            negative <- true
If not positive or not negative:
    in "NIE"; dung

Buoc 3:
Doi cho x va y trong moi vector, roi lap lai Buoc 1 va Buoc 2.
In "TAK".
\end{verbatim}

\subsection{Tiểu kết}

Trong ví dụ trên, mục tiêu ban đầu là tìm nghiệm nguyên không âm của phương
trình. Trực tiếp tìm nghiệm nguyên không âm không dễ, còn tìm một nghiệm
nguyên thì tương đối dễ hơn. Vì vậy trước hết ta phóng đại mục tiêu thành tìm
nghiệm nguyên, rồi từ nghiệm nguyên dựng nghiệm nguyên không âm. Có thể thấy
bài này dùng chuyển hóa mục tiêu trong thiết kế thuật toán; biến bài toán tìm
nghiệm nguyên không âm thành bài toán tìm nghiệm nguyên là chìa khóa.

Trong quá trình giải nghiệm nguyên và nghiệm nguyên không âm, ta cũng không
trực tiếp đưa ra ngay điều kiện cần và đủ. Ta trước hết dựng một điều kiện đủ,
rồi chứng minh tính cần thiết của nó.

\section{Chuyển hóa mục tiêu trong đường lối suy nghĩ}

Vì bài toán phức tạp, suy nghĩ của ta rất dễ trở nên hỗn loạn. Muốn thoát ra
nhưng lại khó tìm điểm đột phá. Khi đó có thể thử chuyển hóa mục tiêu và suy
nghĩ trong mục tiêu đã chuyển hóa. Cách này cho ta một hướng suy nghĩ đúng,
làm mạch lập luận sáng rõ hơn. Ví dụ thường gặp nhất là khi cần xét một bài
toán trong không gian ba chiều, ta có thể hạ nó xuống mặt phẳng hai chiều, rồi
từ mặt phẳng mở rộng trở lại không gian. Ngoài ra còn có cách đặc biệt hóa đồ
thị tổng quát, v.v. Dưới đây thông qua một bài toán để xem sức mạnh của tư
tưởng chuyển hóa mục tiêu trong việc làm rõ suy nghĩ.

\subsection*{Ví dụ 2: Trò chơi bàn cờ}

Cho một bàn cờ. Bàn cờ là một đồ thị vô hướng có \(n\) đỉnh,
\(1\le n\le 100\), và \(m\) cạnh vô hướng. Có hai người chơi, Giáp và Ất. Ban
đầu Giáp đặt một quân cờ lên một đỉnh nào đó; sau đó Ất và Giáp lần lượt di
chuyển quân cờ, Ất đi trước. Mỗi lần chỉ được di chuyển tới một đỉnh kề với
đỉnh hiện tại. Nếu một người đi tới một đỉnh mà quân cờ đã từng đi qua, hoặc
không còn nơi nào để đi, thì người đó thua.

Bây giờ, với một bàn cờ cho trước, cần phán đoán ban đầu Giáp đặt quân cờ ở
những vị trí nào thì có chiến lược thắng. Trong ví dụ ở bản gốc, các đỉnh
\(A,E,F,G\) là các vị trí thắng.

Hình ví dụ trong bản gốc là một đồ thị nhỏ trên các đỉnh
\(A,B,C,D,E,F,G\). Từ hình có thể đọc được phần lập luận cần thiết: nếu Giáp
đặt quân ở \(A\), Ất chỉ có thể đưa quân tới \(B\). Khi đó Giáp đưa quân tới
\(E\). Nếu Ất tiếp tục đi tới \(G\), Giáp đưa quân tới \(F\); nếu Ất đi tới
\(F\), Giáp đưa quân tới \(G\). Trong cả hai trường hợp, Ất đều hết đường đi,
nên Giáp có chiến lược thắng.

\subsection{Thu hẹp phạm vi, thay đổi mục tiêu}

Đây là một bài toán trò chơi đơn giản. Nói nó đơn giản không phải vì bài dễ,
mà vì mô tả bài toán đơn giản và luật chơi cũng đơn giản. Với trò chơi tổng
quát, ta thường dùng cây trò chơi; nhưng với bài này, số trạng thái là cấp số
mũ, nên trực tiếp dùng cây trò chơi khá khó.

Nếu xem bàn cờ là một đồ thị vô hướng, ta có một đồ thị hoàn toàn tổng quát,
không có hạn chế nào, nên khó thảo luận. Vì vậy ta có thể thu hẹp phạm vi đồ
thị bằng cách thêm một số ràng buộc.

Ràng buộc nào? Hãy nhìn một bài tương tự của POI là \textit{GreenGame}. Đó
cũng là một trò chơi trên đồ thị. Hai bài có nhiều điểm giống nhau: đều là trò
chơi, đều diễn ra trên một đồ thị. Nhưng bài này khác \textit{GreenGame} ở
một điểm lớn: đồ thị ở đây là đồ thị tổng quát, còn \textit{GreenGame} là đồ
thị hai phía.

Chính vì vậy, ta thử đơn giản hóa đồ thị thành đồ thị hai phía. Cách này có
lợi gì?

Trước hết, nó bảo đảm khi quân cờ tới một vị trí nào đó thì chắc chắn đến lượt
một người xác định đi. Quan trọng hơn, sự tồn tại của một đỉnh nào đó nhất
định có lợi cho một người và bất lợi cho người còn lại.

Giả sử đồ thị hai phía là
\[
  G=(A,B,E),
\]
trong đó \(A,B\) là hai tập đỉnh và mỗi cạnh trong \(E\) nối một đỉnh của
\(A\) với một đỉnh của \(B\). Không mất tính tổng quát, giả sử ban đầu Giáp
chỉ được đặt quân ở các đỉnh của \(A\). Có thể nói \(A\) thuộc về Giáp và
\(B\) thuộc về Ất. Vì vậy mỗi lần di chuyển là một người chuyển quân từ đỉnh
thuộc về đối phương sang đỉnh thuộc về mình.

Nhìn ở tầm vĩ mô, càng có nhiều ô trong \(A\), Giáp càng có lợi.

\subsection{Liên tục phỏng đoán và chứng minh trên mục tiêu đã chuyển hóa}

Giống các trò chơi thông thường, trước hết ta quy định đỉnh thắng và đỉnh
thua. Giả sử ban đầu quân cờ nằm trên một đỉnh nào đó, có thể thuộc \(A\) hoặc
\(B\); khi ấy người đi trước đã được xác định. Nếu chủ sở hữu của đỉnh đó
\(A\) thuộc Giáp, \(B\) thuộc Ất có chiến lược thắng, ta gọi đó là đỉnh thắng;
ngược lại là đỉnh thua. Từ đây ta nghiên cứu quy luật giữa đỉnh thắng và đỉnh
thua, bắt đầu từ đỉnh thắng và từng bước tìm quan hệ giữa chúng.

\paragraph{Kết luận 3.2.1.}
Hai đỉnh nối bởi cùng một cạnh không thể đều là đỉnh thắng.

Nếu \(p\in A\), \(q\in B\), \((p,q)\in E\), và \(p\) là đỉnh thắng, thì nếu
ban đầu quân ở \(q\), Giáp có thể đưa quân từ \(q\) tới \(p\). Với Ất, các
đỉnh thuộc \(B\) đã đi qua càng ít càng tốt. Khi chưa đi qua \(q\), xuất phát
từ \(p\) Ất đã không có khả năng thắng vì \(p\) là đỉnh thắng của Giáp. Nay
\(q\) đã không thể đi nữa, vẫn xuất phát từ \(p\), Ất càng không thể thắng.
Vì vậy \(q\) tất nhiên là thế thua.

\paragraph{Kết luận 3.2.2.}
Nếu \(p\) là đỉnh thắng và \((p,q)\in E\), thì theo kết luận 3.2.1, \(q\) là
đỉnh thua. Khi đó nhất định tồn tại một đỉnh \(u\), \((u,q)\in E\), \(u\ne p\),
sao cho \(u\) là đỉnh thắng đối với đồ thị
\[
  G-p-q,
\]
tức đồ thị thu được sau khi xóa \(p,q\) và các cạnh liên thuộc với chúng.

Từ nguyên tắc sau một thế thắng nhất định có một thế thua, kết luận này khá dễ
hiểu.

Các kết luận trên cho ta hướng suy nghĩ. Nếu một đỉnh \(p\) trong \(A\) là
đỉnh thắng, thì bất kể bước tiếp theo \(B\) đi tới đỉnh \(q\) nào, \(A\) đều
tìm được một đỉnh \(m(q)\) kề với \(q\) sao cho trong đồ thị đã xóa các đỉnh
đã đi qua, đỉnh ấy là đỉnh thắng.

Đến đây, ta không thể không nghĩ: liệu \(m(q)\) có thể được tính ngay từ đầu
hay không? Nói cách khác, \(m(q)\) có chỉ phụ thuộc vào đồ thị ban đầu không,
bất kể sau này đồ thị thay đổi ra sao? Nếu có, làm thế nào tính \(m(q)\)?

Lần này ta tạm tránh câu hỏi thứ nhất, trực tiếp giả sử câu trả lời là có, rồi
xét cách chơi của Giáp để tìm \(m(q)\).

Trong suy nghĩ của Giáp, với một số đỉnh \(q\) trong \(B\), ta đã tìm được một
đỉnh \(m(q)\). Mỗi khi Ất đưa quân tới \(q\), bước sau Giáp đưa quân tới
\(m(q)\), xét trên đồ thị đã rút gọn. Nếu các \(q\) khác nhau có các
\(m(q)\) khác nhau, cách chơi ấy hiển nhiên khả thi. Nhưng làm sao bảo đảm hai
đỉnh thua khác nhau ứng với hai đỉnh thắng khác nhau? Câu trả lời là: một
ghép cặp trong đồ thị hai phía.

Vậy trong suy nghĩ của Giáp đã có một ghép cặp \(m\), và ghép cặp ấy không
chứa đỉnh ban đầu đặt quân. Giáp vẫn di chuyển theo ký hiệu của \(m\). Nhưng
không phải mọi đỉnh \(q\) trong \(B\) đều có ghép. Nếu cuối cùng đi tới một
đỉnh \(q\) không được ghép thì sao? Từ đỉnh đặt quân ban đầu, đi tới một đỉnh
trong \(B\), rồi tới đỉnh ghép với nó trong \(A\), rồi lại tới một đỉnh trong
\(B\), rồi tới đỉnh ghép với nó trong \(A\), cứ thế tiếp tục, cuối cùng lại
tới một đỉnh \(q\) không được ghép. Từ một đỉnh không ghép trong \(A\), đi xen
kẽ theo cạnh ghép và cạnh không ghép, cuối cùng tới một đỉnh không ghép: đó
chính là một đường tăng.

Vì vậy, ghép cặp \(m\) trong suy nghĩ của Giáp cần bảo đảm không tồn tại đường
tăng. Một ghép cặp không có đường tăng là ghép cặp cực đại về kích thước, tức
ghép cặp lớn nhất.

Như vậy, nếu Giáp đặt quân ở một đỉnh không được ghép trong một ghép cặp lớn
nhất và đi theo cách trên, Giáp nhất định thắng. Nói cách khác, trong ghép cặp
lớn nhất ấy, mọi đỉnh thuộc \(A\) không được ghép đều là đỉnh thắng. Hình
minh họa trong bản gốc tô cạnh của một ghép cặp lớn nhất, đỉnh hiện tại, các
đỉnh đã đi qua và phần cây chiến lược của Giáp; vì màu và mũi tên không khôi
phục được trong trích xuất văn bản, ở đây giữ lại nội dung bằng mô tả trên.

Vì ghép cặp lớn nhất có thể không duy nhất, Giáp có thể tùy ý giả định một
ghép cặp trong đầu. Ta đã tìm được một điều kiện đủ để một đỉnh là đỉnh thắng.

\paragraph{Kết luận 3.2.3.}
Nếu tồn tại một ghép cặp lớn nhất không phủ đỉnh \(p\in A\), thì \(p\) nhất
định là đỉnh thắng.

Vấn đề tiếp theo rất đơn giản: cần chứng minh điều kiện đủ này cũng là cần
thiết.

\paragraph{Kết luận 3.2.4.}
Nếu mọi ghép cặp lớn nhất đều phủ đỉnh \(p\in A\), thì \(p\) nhất định là
đỉnh thua.

\paragraph{Chứng minh.}
Xét từ góc nhìn của Ất. Trong suy nghĩ của Ất cũng có một ghép cặp lớn nhất,
và ghép cặp này nhất định chứa \(p\). Mỗi lần Ất đưa quân tới vị trí được ghép
với vị trí hiện tại. Nếu có một đỉnh không được ghép xuất hiện, thì từ một
đỉnh \(p\) có ghép trong \(A\), ta đi xen kẽ theo cạnh ghép và cạnh không
ghép, cuối cùng tới một đỉnh không được ghép trong \(A\). Nếu đổi trạng thái
ghép và không ghép trên đường này, ta thu được một ghép cặp khác. Số cạnh ghép
không đổi, nên nó vẫn là ghép cặp lớn nhất, nhưng ghép cặp lớn nhất này lại
không chứa \(p\). Điều đó mâu thuẫn với giả thiết mọi ghép cặp lớn nhất đều
phủ \(p\). Vậy trường hợp ấy không thể xảy ra, tức \(p\) nhất định là đỉnh
thua.

Đến đây ta phát hiện câu hỏi đầu tiên ở trên có đáp án đúng, và trường hợp đồ
thị hai phía đã được giải quyết.

\subsection{Quay lại bài toán ban đầu}

Sau khi giải được đồ thị hai phía, ta quay lại đồ thị tổng quát. Lúc này một
đỉnh bất kỳ có thể thuộc về bất cứ người nào; nếu tiếp tục dùng lập luận trên
nguyên xi thì không được. Nhưng dựa trên kết luận của đồ thị hai phía, ta thử
đoán: liệu kết luận 3.2.3 và 3.2.4 có đúng với đồ thị bất kỳ hay không? Không
khó thấy là đúng.

Trong suy nghĩ của Giáp có một ghép cặp lớn nhất. Trước hết Giáp đặt quân ở
một đỉnh không được ghép. Mỗi lần Ất đưa quân tới một đỉnh, Giáp chỉ cần đưa
quân tiếp tới đỉnh ghép với nó. Nếu có một đỉnh không được ghép, ta tìm được
một đường tăng, điều này là không thể.

Về ghép cặp lớn nhất trong đồ thị bất kỳ, có thể tham khảo phần giới thiệu
thuật toán hoa nở trong tài liệu [1].

\subsection{Tiểu kết}

Trong ví dụ trên, đồ thị hai phía không có tác dụng trực tiếp trong thuật toán
cuối cùng, nhưng trong quá trình suy nghĩ nó đem lại rất nhiều thuận tiện.
Mạch suy luận trở nên rõ ràng; các phỏng đoán và chứng minh đều đơn giản hơn.
Có thể nói, khi thu hẹp bài toán về đồ thị hai phía, ta đã giải xong một nửa
bài toán.

Trong đa số trường hợp, không may mắn như ví dụ này: kết luận trên mục tiêu
đã chuyển hóa không nhất thiết chính là kết luận của bài toán ban đầu. Thường
ta còn phải trải qua quá trình phỏng đoán và mở rộng; thậm chí có lúc mục tiêu
đã chuyển hóa khác rất xa bài toán ban đầu, rất khó mở rộng. Nhưng tác dụng
định hướng tư duy của chuyển hóa mục tiêu là điều không thể phủ nhận.

\section{Tổng kết}

Từ những ví dụ trên có thể thấy mục đích của chuyển hóa mục tiêu là làm cho
bài toán đơn giản hơn, bất kể nhu cầu xuất phát từ thuật toán hay từ tư duy.
Nó cung cấp một phương hướng giải quyết tổng quát: từ dễ đến khó, từ đơn giản
đến phức tạp, từ nông đến sâu, từ đặc biệt đến tổng quát. Điều này cũng phù
hợp với hướng vận động của tư duy con người. Vậy làm sao vận dụng được nó một
cách linh hoạt? Không có đáp án rõ ràng, nhưng vài điểm sau rất quan trọng:
\begin{itemize}
  \item Dám phỏng đoán và liên tưởng.
  \item Dám so sánh tương tự và tìm điểm giống nhau.
  \item Dám mở rộng, phát hiện bài toán mới trên nền bài toán đã giải.
\end{itemize}

Cuối cùng, sự thành thạo và linh hoạt vẫn đến từ thực hành. Chỉ khi làm nhiều
và suy nghĩ nhiều, ta mới đạt tới mức vận dụng tự nhiên.

\section*{Lời cảm ơn}

Tác giả cảm ơn Jin Kai, Hu Weidong, He Lin và thầy Xiang Qizhong vì sự giúp
đỡ nhiệt tình.

\section*{Tài liệu tham khảo}

\begin{enumerate}
  \item Wu Wenhu, Wang Jiande. \textit{Hướng dẫn Olympic Tin học quốc tế và
  toàn quốc cho thanh thiếu niên: thuật toán và lập trình đồ thị}.
  \item Wu Wenhu, Wang Jiande. \textit{Phân tích và thiết kế chương trình cho
  các thuật toán thực dụng}.
\end{enumerate}

\end{document}
